\documentclass[11pt]{amsart}
\usepackage[margin=28mm]{geometry}
\usepackage{amssymb,booktabs}
\usepackage{tikz}
\usetikzlibrary{arrows.meta}
\usepackage[hidelinks]{hyperref}
\newtheorem{theorem}{Theorem}
\newtheorem{proposition}[theorem]{Proposition}
\newtheorem{lemma}[theorem]{Lemma}
\theoremstyle{remark}
\newtheorem*{remark}{Remark}
\newcommand{\symking}{\textsf{\textbf{K}}}
\newcommand{\symqueen}{\textsf{\textbf{Q}}}
\newcommand{\symrook}{\textsf{\textbf{R}}}
\newcommand{\symbishop}{\textsf{\textbf{B}}}
\newcommand{\symknight}{\textsf{\textbf{N}}}

\title{Seven problems from Erich Friedman's \emph{Math Magic}}
\author{Vamshi Jandhyala}
\date{Preprint v2, 9 October 2026 (v1 the same day had Sections 2--5). Not peer reviewed.}

\begin{document}
\begin{abstract}
Erich Friedman's \emph{Math Magic} pages pose a problem each month and keep a list of the questions
that remained open. We settle four of them and make progress on three more. Two exact values for armies of bishops,
$B(3,5)=5$ and $B(3,6)=8$, follow from a one-line lemma and an exhaustive count. Three colours of
king with touch counts $(1,3,6)$, $(1,6,3)$ or $(2,3,4)$ admit no finite arrangement (unsolved
problem 24), by a reduction to finite formulas with checkable refutations. Three of the
6-point capture digraphs of unsolved problem 21 are not realised by any chess position. And one of
the heptominoes of unsolved problem 15 can never be arranged in a square with equal row and column
counts, by an explicit weighting. A cube tiled by one $k$-slab for each $k\le n$ has side at most
$n(n+1)/2$, and none exists for $n=4,5,6$ (unsolved problem 6). One entry of the prime-signature
table is misfiled and the smallest number in a cell marked unknown is found, while the two open
cases have no solution below $10^{17}$ (unsolved problem 28). The optimal tournaments of unsolved
problem 33 are computed exactly, improving two published values. Every computation is reproduced
by programs published with the paper.
\end{abstract}
\maketitle

\section{Introduction}

Since 1998 Erich Friedman's \emph{Math Magic} has posed one problem a month, published the
solutions readers sent in, and collected the questions they could not answer on a list of unsolved
problems~\cite{unsolved}. The problems are recreational in flavour and often easy to state, which
makes the open ones tantalising. This note settles four of them and makes progress on three more. As of 9 October 2026 the list still
shows problems 15, 21 and 24 as open, and the March 2005 page on armies of bishops still asks for the
two values we prove~\cite{bishops}.

The results are of two kinds, and we label them. A \emph{proved} statement has a proof a reader can
follow by hand. A \emph{computer-certified} statement rests on a finite computation: an exhaustive
count, or a solver's refutation of a finite formula. For each certified statement we say what was
computed, why the finite computation implies the statement, and how it was checked independently.
The programs are in the \path{verify} folder that accompanies this paper; each section names the
subfolder.

\begin{center}
\small
\begin{tabular}{@{}llll@{}}
\toprule
Problem & Source & Result & Kind\\
\midrule
Armies of bishops & March 2005 & $B(3,5)=5$, $B(3,6)=8$ & computer-certified\\
Touch cycles (no.~24) & March 2015 & three triples impossible & computer-certified\\
Capture digraphs (no.~21) & October 2013 & diagrams $q$, $r$, $s$ impossible & computer-certified\\
Magic polyomino squares (no.~15) & May 2007 & heptomino $A$ impossible & proved\\
Slab cubes (no.~6) & December 2001 & side $\le n(n+1)/2$; none for $n=4,5,6$ & proved; certified\\
Prime signatures (no.~28) & February 2018 & cell $(4,41)$ misfiled; cell $(4,411)$ & proved\\
Optimal tournaments (no.~33) & February 2020 & exact optima & computer-certified\\
\bottomrule
\end{tabular}
\end{center}

\section{Armies of bishops}\label{sec:bishops}

Place $C$ armies of $k$ bishops each on an $N\times N$ board so that no bishop attacks a bishop of
another army. A bishop attacks along diagonals and is blocked by the first piece in its way, of
either army. Let $B(C,N)$ be the largest such $k$. Friedman's table~\cite{bishops} gives lower bounds
for small $C$ and $N$, and asks whether $B(3,5)=5$ and $B(3,6)=8$. Figure~\ref{fig:bishops} shows
three armies of eight on the $6\times6$ board. The question is whether nine is possible.

\begin{figure}[ht]
\centering
\begin{tikzpicture}[scale=0.62]
\fill[white] (0,0) rectangle (1,1);
\fill[black!8] (0,1) rectangle (1,2);
\fill[white] (0,2) rectangle (1,3);
\fill[black!8] (0,3) rectangle (1,4);
\fill[white] (0,4) rectangle (1,5);
\fill[black!8] (0,5) rectangle (1,6);
\fill[black!8] (1,0) rectangle (2,1);
\fill[white] (1,1) rectangle (2,2);
\fill[black!8] (1,2) rectangle (2,3);
\fill[white] (1,3) rectangle (2,4);
\fill[black!8] (1,4) rectangle (2,5);
\fill[white] (1,5) rectangle (2,6);
\fill[white] (2,0) rectangle (3,1);
\fill[black!8] (2,1) rectangle (3,2);
\fill[white] (2,2) rectangle (3,3);
\fill[black!8] (2,3) rectangle (3,4);
\fill[white] (2,4) rectangle (3,5);
\fill[black!8] (2,5) rectangle (3,6);
\fill[black!8] (3,0) rectangle (4,1);
\fill[white] (3,1) rectangle (4,2);
\fill[black!8] (3,2) rectangle (4,3);
\fill[white] (3,3) rectangle (4,4);
\fill[black!8] (3,4) rectangle (4,5);
\fill[white] (3,5) rectangle (4,6);
\fill[white] (4,0) rectangle (5,1);
\fill[black!8] (4,1) rectangle (5,2);
\fill[white] (4,2) rectangle (5,3);
\fill[black!8] (4,3) rectangle (5,4);
\fill[white] (4,4) rectangle (5,5);
\fill[black!8] (4,5) rectangle (5,6);
\fill[black!8] (5,0) rectangle (6,1);
\fill[white] (5,1) rectangle (6,2);
\fill[black!8] (5,2) rectangle (6,3);
\fill[white] (5,3) rectangle (6,4);
\fill[black!8] (5,4) rectangle (6,5);
\fill[white] (5,5) rectangle (6,6);
\draw (0,0) grid (6,6);
\filldraw[fill=black!35,draw=black,line width=0.6pt] (0.5,0.5) circle (0.3);
\filldraw[fill=white,draw=black,line width=0.6pt] (0.5,1.5) circle (0.3);
\filldraw[fill=black,draw=black,line width=0.6pt] (0.5,4.5) circle (0.3);
\filldraw[fill=black!35,draw=black,line width=0.6pt] (0.5,5.5) circle (0.3);
\filldraw[fill=white,draw=black,line width=0.6pt] (1.5,0.5) circle (0.3);
\filldraw[fill=black!35,draw=black,line width=0.6pt] (1.5,1.5) circle (0.3);
\filldraw[fill=white,draw=black,line width=0.6pt] (1.5,2.5) circle (0.3);
\filldraw[fill=black,draw=black,line width=0.6pt] (1.5,3.5) circle (0.3);
\filldraw[fill=black!35,draw=black,line width=0.6pt] (1.5,4.5) circle (0.3);
\filldraw[fill=black,draw=black,line width=0.6pt] (1.5,5.5) circle (0.3);
\filldraw[fill=white,draw=black,line width=0.6pt] (2.5,1.5) circle (0.3);
\filldraw[fill=black,draw=black,line width=0.6pt] (2.5,4.5) circle (0.3);
\filldraw[fill=black,draw=black,line width=0.6pt] (3.5,1.5) circle (0.3);
\filldraw[fill=white,draw=black,line width=0.6pt] (3.5,4.5) circle (0.3);
\filldraw[fill=black,draw=black,line width=0.6pt] (4.5,0.5) circle (0.3);
\filldraw[fill=black!35,draw=black,line width=0.6pt] (4.5,1.5) circle (0.3);
\filldraw[fill=black,draw=black,line width=0.6pt] (4.5,2.5) circle (0.3);
\filldraw[fill=white,draw=black,line width=0.6pt] (4.5,3.5) circle (0.3);
\filldraw[fill=black!35,draw=black,line width=0.6pt] (4.5,4.5) circle (0.3);
\filldraw[fill=white,draw=black,line width=0.6pt] (4.5,5.5) circle (0.3);
\filldraw[fill=black!35,draw=black,line width=0.6pt] (5.5,0.5) circle (0.3);
\filldraw[fill=black,draw=black,line width=0.6pt] (5.5,1.5) circle (0.3);
\filldraw[fill=white,draw=black,line width=0.6pt] (5.5,4.5) circle (0.3);
\filldraw[fill=black!35,draw=black,line width=0.6pt] (5.5,5.5) circle (0.3);
\end{tikzpicture}

\caption{Three armies of eight bishops (black, white, grey) on the $6\times 6$ board. Every
diagonal that carries bishops carries only one army.}\label{fig:bishops}
\end{figure}

The figure suggests the right way to look at a placement: by diagonals, not by bishops.

\begin{lemma}[proved]\label{lem:diag}
A placement of armies is valid if and only if every diagonal containing bishops contains bishops of
only one army.
\end{lemma}

\begin{proof}
If a diagonal carries bishops of two armies, walk along it: somewhere a bishop of one army is
followed, with nothing in between, by a bishop of the other, and these two attack each other.
Conversely, an attack always runs along a diagonal between two bishops on it, so if every diagonal
carries one army no bishop attacks another army.
\end{proof}

So blocking plays no part, and a placement is determined, up to which cells are used, by a label on
each diagonal: the army it carries, or ``unused''. A bishop of army $i$ can stand on a cell exactly
when both diagonals through the cell carry label $i$. Taking every such cell gives a valid placement
by Lemma~\ref{lem:diag}, and removing bishops keeps a placement valid. Write $u_i(\lambda)$ for the number of cells whose two diagonals both carry label $i$ in the
labelling $\lambda$. Hence
\[
B(C,N)=\max_{\lambda}\ \min_{1\le i\le C}\ u_i(\lambda).
\]
The right side is the size of the smallest army when every usable cell is filled, maximised over all
ways of labelling the diagonals. A light square and a dark square never share a diagonal, so the
light and dark diagonals can be labelled separately: for each colour we list the vectors of army sizes
that its labellings achieve, then combine one vector from each colour.

\begin{theorem}[computer-certified]\label{thm:bishops}
$B(3,5)=5$ and $B(3,6)=8$.
\end{theorem}

\begin{proof}[Method]
The lower bounds are the placements of Figure~\ref{fig:bishops} and of its $5\times5$ analogue in
\path{verify/bishops/positions.json}, checked by a move-by-move attack simulator. For the upper
bounds, with $C=3$ each diagonal takes one of four labels. On the $5\times5$ board the light squares
lie on ten diagonals and the dark squares on eight, giving $4^{10}+4^{8}$ labellings; on the
$6\times6$ board each colour has eleven, giving $2\cdot4^{11}$. The program
\path{verify/bishops/check.py} lists them all and evaluates the formula above. The maxima are $5$ and
$8$. The same program reproduces $B(3,3)=2$ and $B(3,4)=4$, and with two armies the values
$B(2,5)=12$ and $B(2,6)=18$ in Friedman's table. In \path{mutants.py} a recoloured bishop is caught
by the simulator, and a mutant that lets a bishop stand on a cell when only one of its diagonals
matches its army gives a larger value for $B(3,5)$, so the count does depend on the lemma.
\end{proof}

\section{Touch cycles (unsolved problem 24)}\label{sec:touch}

Place kings of three colours $A$, $B$, $C$ on the cells of the plane, at most one per cell. Two cells
touch when they are king moves apart, that is, when their coordinates differ by at most one in each
direction. Friedman~\cite{touch} asks for which triples $(a,b,c)$ there is a nonempty finite
arrangement in which every $A$ touches exactly $a$ kings of colour $B$, every $B$ touches exactly $b$
of colour $C$, and every $C$ touches exactly $c$ of colour $A$. Figure~\ref{fig:touch} shows one for
$(1,1,6)$. The unsolved list gives $(1,3,6)$, $(1,6,3)$ and $(2,3,4)$ as unknown and conjectures that
they do not exist.

\begin{figure}[ht]
\centering
\begin{tikzpicture}[scale=0.55]
\draw[black!25] (-0.5,0.5) grid (4.5,6.5);
\filldraw[fill=black,draw=black,line width=0.6pt] (0.5,1.5) circle (0.32);
\filldraw[fill=white,draw=black,line width=0.6pt] (0.5,2.5) circle (0.32);
\filldraw[fill=white,draw=black,line width=0.6pt] (0.5,3.5) circle (0.32);
\filldraw[fill=black,draw=black,line width=0.6pt] (0.5,4.5) circle (0.32);
\filldraw[fill=white,draw=black,line width=0.6pt] (0.5,5.5) circle (0.32);
\filldraw[fill=white,draw=black,line width=0.6pt] (1.5,1.5) circle (0.32);
\filldraw[fill=black!35,draw=black,line width=0.6pt] (1.5,2.5) circle (0.32);
\filldraw[fill=white,draw=black,line width=0.6pt] (1.5,3.5) circle (0.32);
\filldraw[fill=black!35,draw=black,line width=0.6pt] (1.5,4.5) circle (0.32);
\filldraw[fill=white,draw=black,line width=0.6pt] (1.5,5.5) circle (0.32);
\filldraw[fill=white,draw=black,line width=0.6pt] (2.5,1.5) circle (0.32);
\filldraw[fill=black!35,draw=black,line width=0.6pt] (2.5,2.5) circle (0.32);
\filldraw[fill=white,draw=black,line width=0.6pt] (2.5,3.5) circle (0.32);
\filldraw[fill=black!35,draw=black,line width=0.6pt] (2.5,4.5) circle (0.32);
\filldraw[fill=white,draw=black,line width=0.6pt] (2.5,5.5) circle (0.32);
\filldraw[fill=white,draw=black,line width=0.6pt] (3.5,1.5) circle (0.32);
\filldraw[fill=black,draw=black,line width=0.6pt] (3.5,2.5) circle (0.32);
\filldraw[fill=white,draw=black,line width=0.6pt] (3.5,3.5) circle (0.32);
\filldraw[fill=white,draw=black,line width=0.6pt] (3.5,4.5) circle (0.32);
\filldraw[fill=black,draw=black,line width=0.6pt] (3.5,5.5) circle (0.32);
\end{tikzpicture}

\caption{A finite arrangement for $(a,b,c)=(1,1,6)$: every $A$ (white) touches one $B$ (black),
every $B$ touches one $C$ (grey), and every $C$ touches six $A$'s.}\label{fig:touch}
\end{figure}

\begin{theorem}[computer-certified]\label{thm:touch}
There is no nonempty finite arrangement for $(a,b,c)=(1,3,6)$, $(1,6,3)$ or $(2,3,4)$. Nor is there
one for $(1,4,5)$.
\end{theorem}

\begin{proof}[Method]
The idea is to look at an extreme king, where the arrangement must stop. Fix a direction $n$, namely
$n=(0,1)$ for the first, second and fourth triples and $n=(1,1)$ for $(2,3,4)$. Given a finite
arrangement, choose the occupied cell $p$ that maximises $n\cdot p$, breaking ties by the larger
first coordinate, and translate so that $p$ is the origin. Then every cell $q$ with $n\cdot q>0$, or
with $n\cdot q=0$ and positive first coordinate, is empty.

Now restrict attention to the window $[-r,r]^2$. Every cell strictly inside it, with both
coordinates of absolute value less than $r$, has all eight neighbours in the window, so the counting
rule at that cell refers only to cells of the window. The restriction of the arrangement is therefore
a solution of a finite problem: colour each window cell $A$, $B$, $C$ or empty, with the origin
occupied, the half-plane beyond it empty, and the counting rule imposed at every occupied interior
cell. That problem is unsatisfiable for $r=6$, $6$, $8$ and $4$ for the four triples in the order
listed. So no finite arrangement exists.

Each finite problem is written as a formula in conjunctive normal form, and the folder
\path{verify/touch-cycles/certificates} holds the formula with a DRUP refutation found by a SAT
solver. The checker \path{verify_drup.py}, which uses only the Python standard library, confirms
all four refutations in about thirteen seconds. Trust then moves to the encoding, and that was
checked separately: \path{touch_extreme.py} is an independently written encoding of the same window
problem for the Z3 solver. It reports the four problems unsatisfiable at the same radii, and
satisfiable for the controls $(1,2,3)$ and $(2,2,2)$, which do have finite arrangements.
\end{proof}

For $(1,4,5)$, which is not on the list, there is also a proof by hand
(\path{HUMAN-PROOF-145.md}). For the three listed triples finiteness is essential: each has
a doubly periodic arrangement on the whole plane, with periods $6\times6$, $8\times8$ and $6\times6$
(\path{touch_torus.py}). A proof of Theorem~\ref{thm:touch} must therefore use the boundary of the
arrangement, as the extreme-king argument does.

\section{Capture digraphs (unsolved problem 21)}\label{sec:capture}

Put chess pieces (kings, queens, rooks, bishops and knights, any number of each, colours ignored) on
a board, and draw an arrow from piece $X$ to piece $Y$ when $X$ attacks $Y$; sliding pieces are
blocked by the first piece in their way. Friedman~\cite{capture} asks which regular digraphs arise
this way. For six vertices with two arrows in and two out at every vertex there are 23 digraphs up to
isomorphism, which his page labels $a$ to $w$. Readers found positions for eighteen of them,
DeVincentis showed that $f$ cannot be realised, and the author found a position for $u$ on a
$3\times4$ board (Figure~\ref{fig:capture}, left). The diagrams $q$, $r$ and $s$ remained open, and
the unsolved list conjectures that they are not capture graphs.

\begin{figure}[ht]
\centering
\begin{tikzpicture}[scale=0.9]
\fill[white] (0,0) rectangle (1,1);
\fill[black!8] (0,1) rectangle (1,2);
\fill[white] (0,2) rectangle (1,3);
\fill[black!8] (0,3) rectangle (1,4);
\fill[black!8] (1,0) rectangle (2,1);
\fill[white] (1,1) rectangle (2,2);
\fill[black!8] (1,2) rectangle (2,3);
\fill[white] (1,3) rectangle (2,4);
\fill[white] (2,0) rectangle (3,1);
\fill[black!8] (2,1) rectangle (3,2);
\fill[white] (2,2) rectangle (3,3);
\fill[black!8] (2,3) rectangle (3,4);
\draw (0,0) grid (3,4);
\draw[-{Stealth[length=5pt]},black!60,line width=0.7pt,shorten >=9pt,shorten <=9pt,bend left=12] (2.5,3.5) to (0.5,2.5);
\draw[-{Stealth[length=5pt]},black!60,line width=0.7pt,shorten >=9pt,shorten <=9pt,bend left=12] (2.5,3.5) to (1.5,1.5);
\draw[-{Stealth[length=5pt]},black!60,line width=0.7pt,shorten >=9pt,shorten <=9pt,bend left=12] (0.5,2.5) to (2.5,2.5);
\draw[-{Stealth[length=5pt]},black!60,line width=0.7pt,shorten >=9pt,shorten <=9pt,bend left=12] (0.5,2.5) to (0.5,1.5);
\draw[-{Stealth[length=5pt]},black!60,line width=0.7pt,shorten >=9pt,shorten <=9pt,bend left=12] (2.5,2.5) to (2.5,3.5);
\draw[-{Stealth[length=5pt]},black!60,line width=0.7pt,shorten >=9pt,shorten <=9pt,bend left=12] (2.5,2.5) to (1.5,1.5);
\draw[-{Stealth[length=5pt]},black!60,line width=0.7pt,shorten >=9pt,shorten <=9pt,bend left=12] (0.5,1.5) to (2.5,3.5);
\draw[-{Stealth[length=5pt]},black!60,line width=0.7pt,shorten >=9pt,shorten <=9pt,bend left=12] (0.5,1.5) to (1.5,0.5);
\draw[-{Stealth[length=5pt]},black!60,line width=0.7pt,shorten >=9pt,shorten <=9pt,bend left=12] (1.5,1.5) to (0.5,1.5);
\draw[-{Stealth[length=5pt]},black!60,line width=0.7pt,shorten >=9pt,shorten <=9pt,bend left=12] (1.5,1.5) to (1.5,0.5);
\draw[-{Stealth[length=5pt]},black!60,line width=0.7pt,shorten >=9pt,shorten <=9pt,bend left=12] (1.5,0.5) to (0.5,2.5);
\draw[-{Stealth[length=5pt]},black!60,line width=0.7pt,shorten >=9pt,shorten <=9pt,bend left=12] (1.5,0.5) to (2.5,2.5);
\node[font=\large] at (2.5,3.5) {\symknight};
\node[font=\large] at (0.5,2.5) {\symrook};
\node[font=\large] at (2.5,2.5) {\symking};
\node[font=\large] at (0.5,1.5) {\symbishop};
\node[font=\large] at (1.5,1.5) {\symrook};
\node[font=\large] at (1.5,0.5) {\symknight};
\end{tikzpicture}
\qquad\qquad
\begin{tikzpicture}[scale=1.25]
\node[circle,fill=black,inner sep=2.2pt] (v0) at (90:1.25) {};
\node[circle,fill=black,inner sep=2.2pt] (v1) at (150:1.25) {};
\node[circle,fill=black,inner sep=2.2pt] (v2) at (210:1.25) {};
\node[circle,fill=black,inner sep=2.2pt] (v3) at (270:1.25) {};
\node[circle,fill=black,inner sep=2.2pt] (v4) at (330:1.25) {};
\node[circle,fill=black,inner sep=2.2pt] (v5) at (390:1.25) {};
\draw[-{Stealth[length=6pt]},line width=0.7pt,shorten >=3pt,shorten <=3pt,bend left=14] (v0) to (v1);
\draw[-{Stealth[length=6pt]},line width=0.7pt,shorten >=3pt,shorten <=3pt] (v0) to (v2);
\draw[-{Stealth[length=6pt]},line width=0.7pt,shorten >=3pt,shorten <=3pt,bend left=14] (v1) to (v0);
\draw[-{Stealth[length=6pt]},line width=0.7pt,shorten >=3pt,shorten <=3pt] (v1) to (v2);
\draw[-{Stealth[length=6pt]},line width=0.7pt,shorten >=3pt,shorten <=3pt] (v2) to (v3);
\draw[-{Stealth[length=6pt]},line width=0.7pt,shorten >=3pt,shorten <=3pt] (v2) to (v4);
\draw[-{Stealth[length=6pt]},line width=0.7pt,shorten >=3pt,shorten <=3pt] (v3) to (v0);
\draw[-{Stealth[length=6pt]},line width=0.7pt,shorten >=3pt,shorten <=3pt] (v3) to (v5);
\draw[-{Stealth[length=6pt]},line width=0.7pt,shorten >=3pt,shorten <=3pt] (v4) to (v3);
\draw[-{Stealth[length=6pt]},line width=0.7pt,shorten >=3pt,shorten <=3pt,bend left=14] (v4) to (v5);
\draw[-{Stealth[length=6pt]},line width=0.7pt,shorten >=3pt,shorten <=3pt] (v5) to (v1);
\draw[-{Stealth[length=6pt]},line width=0.7pt,shorten >=3pt,shorten <=3pt,bend left=14] (v5) to (v4);
\end{tikzpicture}

\caption{Left: a position on a $3\times4$ board whose capture graph is the 2-regular digraph $u$;
each arrow is an attack. Right: the digraph $q$, which no position realises.}\label{fig:capture}
\end{figure}

\begin{theorem}[computer-certified]\label{thm:capture}
The digraphs $q$, $r$ and $s$ are not the capture graph of any position, on a board of any size.
\end{theorem}

\begin{proof}[Method]
A position on any finite board is also a position on the infinite board $\mathbb Z^2$ with the same
attacks, so it is enough to rule out $\mathbb Z^2$. The program \path{src/digraphs.py} lists the
2-regular digraphs on six vertices up to isomorphism and finds 23, matching Friedman's count. To tie
the solver's numbering to his letters without reading arrows off pictures, the eighteen published
positions were transcribed (\path{out/published_solutions.txt}) and their capture graphs
computed; they cover eighteen distinct classes, and the five left over are exactly the five diagrams
$f$, $q$, $r$, $s$, $u$ without a published position.

Two models then decide every class. The first (\path{src/realize.py}) gives each piece unbounded
integer coordinates and asks Z3, in linear integer arithmetic, for a position with exactly the
required arrows; it enumerates the assignments of piece types to vertices, discarding only
assignments that cannot work for a stated reason (two pieces of the same type attack each other
symmetrically, for example). The second (\path{verify/second.py}) is written independently for
the cvc5 solver, keeps the piece types as solver variables, and does no pruning. Both models find
positions for 19 classes and prove the other four unrealisable: $f$, $q$, $r$ and $s$. Every
position the solvers return is re-checked by a square-by-square simulator (\path{verify/check.py}),
and the encodings are tested against that simulator on random positions, with deliberate mutants
caught (\path{verify/test_*.py}).
\end{proof}

The result for $f$ agrees with DeVincentis. A proof of Theorem~\ref{thm:capture} by hand would still
be welcome: after the pruning there are only 240, 60 and 60 type assignments to rule out for $q$,
$r$ and $s$.

\section{A heptomino that never balances (unsolved problem 15)}\label{sec:heptomino}

Friedman's May 2007 problem~\cite{polyomino} asks how few copies of a given polyomino can be placed,
without overlapping, in a square so that every row and every column contains the same positive number
of covered cells. Rotations and reflections are allowed. For most small polyominoes some arrangement
works. Unsolved problem 15 lists heptominoes for which none is known and conjectures that none exists.
We settle the first of them, the shape $A$ of Figure~\ref{fig:heptomino}.

\begin{figure}[ht]
\centering
\begin{tikzpicture}[scale=0.6]
\filldraw[fill=black!18,draw=black] (1,4) rectangle (2,5);
\filldraw[fill=black!18,draw=black] (2,4) rectangle (3,5);
\node[font=\small] at (3.6,4.5) {2};
\filldraw[fill=black!18,draw=black] (1,3) rectangle (2,4);
\node[font=\small] at (3.6,3.5) {1};
\filldraw[fill=black!18,draw=black] (1,2) rectangle (2,3);
\node[font=\small] at (3.6,2.5) {1};
\filldraw[fill=black!18,draw=black] (1,1) rectangle (2,2);
\node[font=\small] at (3.6,1.5) {1};
\filldraw[fill=black!18,draw=black] (0,0) rectangle (1,1);
\filldraw[fill=black!18,draw=black] (1,0) rectangle (2,1);
\node[font=\small] at (3.6,0.5) {2};
\node[font=\small] at (0.5,-0.45) {1};
\node[font=\small] at (1.5,-0.45) {5};
\node[font=\small] at (2.5,-0.45) {1};
\end{tikzpicture}

\caption{The heptomino $A$ with its row counts $P=(2,1,1,1,2)$ and column counts $Q=(1,5,1)$.}
\label{fig:heptomino}
\end{figure}

\begin{theorem}[proved]\label{thm:heptomino}
No arrangement of one or more copies of $A$ in an $n\times n$ square has the same positive number of
covered cells in every row and every column.
\end{theorem}

\begin{proof}
Give each column $x$ a weight $w_x$ and each row $y$ the same weight $w_y$, and give the cell $(x,y)$
the weight $w_x+w_y$. The weight of a set of cells is then the sum, over columns, of the column's
weight times the number of cells in it, plus the same sum over rows. For a balanced arrangement with
$c$ covered cells in every row and column this is $2c\sum_x w_x$.

Now weigh a single copy. The row counts of $A$ are $P=(2,1,1,1,2)$ and its column counts are
$Q=(1,5,1)$. Both are palindromes, so in every rotation and reflection one direction sees $P$ across
five consecutive lines and the other sees $Q$ across three. A copy therefore weighs
$P\cdot(w_i,\dots,w_{i+4})+Q\cdot(w_j,\dots,w_{j+2})$ for some $i$ and $j$. Call these the $P$-windows
and $Q$-windows of $w$. If every $P$-window is at least $-1$ and every $Q$-window is at least $1$, each
copy weighs at least $0$, and so does any arrangement. If moreover $\sum w<0$, a balanced arrangement
would weigh $2c\sum w<0$. It remains to find such weights for every $n\ge5$; a square of side less than
five cannot hold $A$ at all.

For $5\le n\le 9$ exact rational weights are listed in \path{verify/heptomino/exact-weights-A.jsonl}.
For larger $n$ write $n=5K+m$ with $K\ge1$ and $5\le m\le 9$, and build $w$ from $K$ blocks of length
five followed by a core:
\[
w=\bigl(a_KZ+d,\ \dots,\ a_1Z+d,\ C_m\bigr),\quad Z=(-5,1,0,-1,5),\quad
d=\bigl(-\tfrac12,\tfrac12,0,\tfrac12,-\tfrac12\bigr),
\]
with $a_1=\tfrac72$, $a_{k+1}=5a_k$, and fixed cores $C_m$ of negative sum given in
\path{shapeA_family.json}. Each block sums to zero, since $Z$ and $d$ do, so $\sum w=\sum C_m<0$. A
window lying inside one block does not depend on its scale: $P\cdot Z=-10+1+0-1+10=0$, and the
three $Q$-windows of $Z$ are $-5+5+0$, $1+0-1$ and $0-5+5$, all zero. Such windows meet the bounds
through $d$. A window across the boundary between blocks of scales $5s$ and $s$ has
the form $\alpha s+\delta$ with $\alpha>0$, and it meets its bound at $s=a_1$, hence at every larger
scale. The windows that touch the core take finitely many values and are checked directly.
\end{proof}

The checker \path{verify/heptomino/check_shapeA.py} uses exact rational arithmetic and only the
standard library. It verifies the window bounds and the negative sum explicitly for every side from 5
to 400, checks the induction step symbolically, and rejects four mutants: a wrong core, the ratio $4$
in place of $5$, no offset $d$, and a wrong profile. The other heptominoes on the list remain open.
For them the weighting method fails: a fractional relaxation of the placement problem is feasible on
some larger squares, so no weights of this kind can exist there.


\section{Slab cubes (unsolved problem 6)}\label{sec:slabs}

A $k$-slab is a $k\times ik\times jk$ box for some positive integers $i,j$, in any orientation. Friedman
asks which cubes can be tiled by one $k$-slab for each $k=1,\dots,n$~\cite{slabs}; the only known ones are
the unit cube ($n=1$) and the $6\times6\times6$ cube ($n=3$), and he conjectures there are infinitely
many. For example the $6^3$ cube is a $1\times6\times6$, a $2\times6\times6$ and a $3\times6\times6$
slab side by side.

\begin{proposition}[proved]
A cube tiled by one $k$-slab for each $k\le n$ has side $s\le n(n+1)/2$, with equality only for $n=1$
and $n=3$.
\end{proposition}

\begin{proof}
A $k$-slab inside the cube has two sides at most $s$, so its volume is at most $ks^2$; hence
$s^3\le s^2(1+\dots+n)$. Equality forces every slab to be a full $k\times s\times s$ layer, so $k$
divides $s=n(n+1)/2$ for every $k\le n$. For $n=2$, $2\nmid3$. For $n\ge4$, $n-1$ would divide
$n(n+1)$, which is $2$ modulo $n-1$, so $n-1\mid2$, impossible.
\end{proof}

\begin{proposition}[computer-certified]
No cube is tiled by one $k$-slab for each $k\le n$ when $n=4$, $5$ or $6$.
\end{proposition}

\noindent By the volume bound and the first proposition, only sides with $\sum k^3\le s^3$ and
$s\le n(n+1)/2$ need checking: $5\le s\le10$, $7\le s\le15$ and $8\le s\le 21$. Two programs decide
every case, in different ways (\path{verify/slab-cubes}). The first lists every choice of slab sizes
whose volumes add up to $s^3$ and refutes each with an exact-cover model in a constraint solver
(\path{slabs.py}, \path{slabs_cube.py}; for $n=6$, $s=12$ there are 672 choices). The second uses no
solver: it fills the cube in lexicographic order, where the first empty cell must be the corner of the
slab covering it, tries every unused slab, size and orientation there, and prunes a branch unless the
empty volume is still a sum of volumes of the unused slabs (\path{slab_search.py}). Both find tilings of
the $6^3$ cube and of Friedman's smallest boxes $2\times2\times3$, $3\times5\times6$ and
$4\times6\times7$, and both report no tiling for every side above. A weakened search that allows any
$k\times i\times j$ box does tile the cubes of sides 5 and 6 with $n=4$, so the negative answers are not
vacuous (\path{slab_mutant.py}). The search for $n=7$ ($10\le s\le
28$) has not been run.

\section{Prime signatures (unsolved problem 28)}\label{sec:primesig}

The prime signature of $n$ is the multiset of exponents in its factorisation, written as a string:
$12=2^2\cdot3$ has signature $21$. Friedman tabulates, for signatures $A$ and $B$ of total degree at most
6, the smallest $n$ of signature $A$ with $n+1$ of signature $B$~\cite{primesig}. The cases $(A,B)=(22,32)$
and $(32,32)$ are listed as open, with the conjecture that they have no solutions.

\begin{proposition}[proved]
The table's entry $n=1022303^4$ in cell $(A,B)=(4,41)$ is misfiled: $n+1=2\cdot17^4\cdot
6538721797256344321$ with the last factor prime, so $n+1$ has signature $411$. Cell $(4,41)$ is empty,
and $1022303^4$ is the smallest $n$ in cell $(4,411)$, which the table marks unknown.
\end{proposition}

\begin{proof}
For odd $p$, $p^4+1\equiv2\pmod{16}$, so $2$ divides $p^4+1$ exactly once, and $p=2$ gives $17$.
Signature $41$ would force $p^4+1=2q^4$, that is $x^2+1=2y^4$ with $x=p^2$; its only positive solutions
are $x=1$ and $x=239$~\cite{ljunggren}, and neither is the square of a prime. Signature $411$ forces
$p^4+1=2q^4r$ with $q\ne r$ odd primes, and $q^4\mid p^4+1$ requires $q\equiv1\pmod 8$. For every such
$q<1022303$ the four roots of $x^4\equiv-1\pmod{q^4}$ were lifted from roots modulo $q$, and every prime
$p<1022303$ in those classes was tested: none gives a prime cofactor $r$
(\path{verify/prime-signatures/primesig_4_411.py}, seconds). The primality of the 19-digit factor is
certified by a deterministic test valid below $2^{64}$.
\end{proof}

\noindent\emph{Evidence for the open cases.} No $n\le10^{17}$ has signature $22$ with $n+1$ of
signature $32$, or signature $32$ with $n+1$ of signature $32$ (\path{primesig_pairs.py},
\path{data/primesig_pairs_17.txt}). Every $p^aq^b\le10^{17}$ of the two shapes is listed, using primes up
to $\sqrt{10^{17}/4}$ so that the numbers $4q^2$ are included; the lists agree with brute-force
factorisation of every $n\le10^7$ (\path{test_primesig_pairs.py}), which also finds the control pair
$675=3^35^2$, $676=2^213^2$. A first run took primes only to $\sqrt{X/8}$ and so missed some numbers $4q^2$;
the test fails on that version, and the corrected run still finds no pair. All 650 published values were
checked by factorisation (649 are correct, the exception being the entry above), and every value below
$3\times10^6$ was confirmed smallest by a sieve (\path{primesig_table.py}).

\section{Optimal tournaments (unsolved problem 33)}\label{sec:tournament}

$K$ players are in a uniformly random order of ability, and the better player wins each game with
probability $2/3$. $N$ games are scheduled adaptively, each chosen after seeing the earlier results, and
then one guesses either the whole order (``rank'') or the best player (``winner'')~\cite{tournament}.
Friedman tabulates the best probabilities readers found. For example with two players and one game,
guessing that the winner is the better player is right with probability $2/3$; the tables record the
best such probabilities for each $K$ and $N$.

\begin{proposition}[computer-certified]
The optimal probabilities that differ from the published table are:
\begin{itemize}
\item rank, $K=3$, $N=7$: $2992/6561$ (published $2944/6561$); rank, $K=4$, $N=7$: $1216/6561$
(published $1136/6561$);
\item rank, $K=5$: for $N=2$ the optimum is $2/135$ (the table prints $2/125$ beside the correct
$1.48\%$); for $N=5$ it is $128/3645$, which the table's own schedule attains (published $16/405$);
for $N=6$ it is $512/10935$ (published $64/1215$, which no schedule attains).
\end{itemize}
Every other entry with $K\le5$, $N\le7$ or $K=6$, $N\le6$ in the rank table, and with $K\le4$, $N\le12$,
$K=5$, $N\le7$, $K=6$, $N\le6$ or $K=7$, $N\le4$ in the winner table, agrees with the published value,
which is therefore optimal. In particular Knop's $1984/32805<1/13$ for rank, $K=5$, $N=7$ is optimal,
consistent with Kopp's disproof of the conjecture that $1/13$ is attainable.
\end{proposition}

\noindent The likelihood of a set $H$ of results under an order $s$ is $2^{a(s,H)}/3^{|H|}$, where $a$
counts the results consistent with $s$. So the optimum is an exact recursion over multisets of
results: choose the next pair to maximise the expected value, and at the end guess the most likely
order (\path{verify/tournaments/tournament.py}). A separate program, which carries the exact
probability of the results under every order instead of the power-of-two count
(\path{tournament_check.py}), agrees at $(K,N)=(3,7)$, $(4,7)$ and
$(5,5)$; the extracted optimal schedule for $K=3$, $N=7$ (\path{data/tournament_rank_3_7.json}) re-scores
to $2992/6561$, and simulating it agrees within sampling error. \path{tournament_table.py} compares
every optimum with the published tables.

\section*{Verification and use of AI tools}

The searches, certificates and parts of the proofs were produced with AI coding assistants, Claude
(Anthropic) and Codex (OpenAI), and then checked as each section describes. In particular the
extreme-king reduction and DRUP certificates for touch cycles, and the finite weights for the
heptomino on small squares, were first found by Codex; the slab, prime-signature and tournament
computations of Sections 6--8 were written with Claude. The author takes responsibility for every
statement. All programs and data are in the folder that accompanies this paper; each section names
its subfolder, and \path{README.md} gives the commands and running times.

\begin{thebibliography}{9}
\bibitem{unsolved} E.~Friedman, \emph{Math Magic: unsolved problems},
\url{https://erich-friedman.github.io/mathmagic/unsolved.html}, accessed 9 October 2026.
\bibitem{bishops} E.~Friedman, \emph{Math Magic}, problem of the month, March 2005 (armies of bishops),
\url{https://erich-friedman.github.io/mathmagic/0305.html}, accessed 9 October 2026.
\bibitem{polyomino} E.~Friedman, \emph{Math Magic}, problem of the month, May 2007 (polyominoes in
squares), \url{https://erich-friedman.github.io/mathmagic/0507.html}, accessed 9 October 2026.
\bibitem{capture} E.~Friedman, \emph{Math Magic}, problem of the month, October 2013 (regular
digraphs of chess pieces), \url{https://erich-friedman.github.io/mathmagic/1013.html}, accessed
9 October 2026.
\bibitem{touch} E.~Friedman, \emph{Math Magic}, problem of the month, March 2015 (touching letters),
\url{https://erich-friedman.github.io/mathmagic/0315.html}, accessed 9 October 2026.
\bibitem{slabs} E.~Friedman, \emph{Math Magic}, problem of the month, December 2001 (slabs),
\url{https://erich-friedman.github.io/mathmagic/1201.html}, accessed 9 October 2026.
\bibitem{primesig} E.~Friedman, \emph{Math Magic}, problem of the month, February 2018 (prime
signatures), \url{https://erich-friedman.github.io/mathmagic/0218.html}, accessed 9 October 2026.
\bibitem{tournament} E.~Friedman, \emph{Math Magic}, problem of the month, February 2020 (tournaments),
\url{https://erich-friedman.github.io/mathmagic/0220.html}, accessed 9 October 2026.
\bibitem{ljunggren} W.~Ljunggren, Zur Theorie der Gleichung $x^2+1=Dy^4$, \emph{Avh. Norske Vid. Akad.
Oslo I. Mat.-Naturv. Kl.} 1942, no.~5, 27~pp.
\end{thebibliography}

\end{document}
